\chapter{MATERIAIS E MÉTODOS}
\label{ch:metodos}

% TODO: descrever as seis rodadas finais com base nas configurações e nos resultados salvos.
\section{DADOS}
% Dois corpora de 50 entradas: descrições políticas em inglês e legendas jornalísticas em português.
% Documentar seleção, limpeza, ordem, identificadores, hashes, origem e licença; não são traduções.

\section{MODELOS, AGENTES E PERSPECTIVAS}
% Separar os três geradores do gpt-6-luna usado na extração, avaliação estrutural e VAD LLM.
% Documentar parâmetros efetivamente enviados: temperatura 0,8 para Grok/Gemini, omitida para GPT-6.

\section{GERAÇÃO DAS CADEIAS DE REESCRITA}
% Oito trajetórias de quatro posições em quatro árvores: NNNN, CCCC, PPPP, DDDD, CCPP, PPCC, DDNN, NNDD.
% Explicar modo interpretativo, NeutralRelay e prefixos compartilhados: 1.200 nós por rodada.
% As 1.600 posições projetadas não são gerações independentes; distinguir falhas, retomadas e cache.

\section{EXTRAÇÃO ESTRUTURAL}
% Registrar versões e ausência do título separado; narrative_frame não tem peso próprio no índice.

\section{STRUCTURED TOPIC DRIFT INDEX}
% Inglês: Dual (Cluster/LLM); português: LLM. Distinguir conteúdo, contradição interna e VAD.
% Três avaliações estruturais por par, com média de STDIs completos e registro de atalhos/ausências.
% Repetição do avaliador não repete geração, extração ou VAD; tentativa após falha e cache são distintos.
% Separar distância ao original, mudança contra o pai efetivo e soma cumulativa, que pode exceder 1.

\section{ANÁLISE DE VALÊNCIA, ATIVAÇÃO E DOMINÂNCIA}
% Inglês: BERT/LLM em VAD Dual; português: LLM. Dual promedia derivas normalizadas, não vetores brutos.
% NRC e MEmoLon são alternativas exploratórias, sem adoção no cálculo das rodadas finais.

\section{PROCEDIMENTOS DE AVALIAÇÃO E ANÁLISE}
% Usar comparações quantitativas, casos controlados e inspeções qualitativas existentes, sem notas humanas esperadas.
% Comparar trajetórias e geradores de forma pareada por entrada dentro de cada corpus, respeitando dependências.
% O ramo LLM/VAD LLM salvo em inglês aproxima o procedimento português, sem isolar o efeito do idioma.
% Não presumir execução de sensibilidades, testes ou intervalos sem localizar o artefato correspondente.
